% GENERATED FILE. Edit canonical lessons, then run npm run generate:books.
% canonical-source-hash: fnv1a64:70660ad5822b59b1
% canonical-lessons: RU-C209-fotografirovat, RU-C209-sobirat, RU-C209-obeshchat, RU-C209-shutit, RU-C209-zhenitsya

\chapter{To photograph, To gather, To promise, To joke, To get married}
\label{ch:ru-to-photograph-to-gather-to-promise-to-joke-to-get-married}
\hlchaptermodality{209}

\begin{chapteropening}
\textbf{By the end of this chapter:} \emph{I can say to photograph, to gather, to promise, to joke and to get married.}

Complete the last lesson of chapter 209: I can say to photograph, to gather, to promise, to joke and to get married.
\end{chapteropening}

\section[\texorpdfstring{\ru{фотографировать} (fotografírovat)}{фотографировать (fotografírovat)}]{\ru{фотографировать} (fotografírovat) — to photograph}

\label{lesson:RU-C209-fotografirovat}



\begin{quote}
Before the new one: say the Russian for to sniff, then the Russian for to touch.
\end{quote}

\subsection*{You'll want to know: \ru{фотографировать}}
\textbf{\ru{фотографировать}} (\emph{fotografírovat}) — \textquotedblleft{}to photograph\textquotedblright{}.

\begin{grammarlens}[title={What you've built}]
One more word for this chapter.
\end{grammarlens}

\subsection*{Your turn}
\begin{itemize}
\raggedright
  \item \emph{Say it:} \emph{fotografírovat}
  \item \emph{Say it:} \emph{fotografírovat}, once more
  \item \emph{Say it:} say \emph{fotografírovat}
  \item \emph{Recall:} say the Russian for to sniff, then the Russian for to touch, then say \emph{fotografírovat} again
\end{itemize}

\begin{tcolorbox}[breakable,colback=teal!4,colframe=teal!35!black,title={Before you move on}]
What does \ru{фотографировать} mean? (\textquotedblleft{}To photograph\textquotedblright{}.) Say it once more.
\end{tcolorbox}
\section[\texorpdfstring{\ru{собирать} (sobirát)}{собирать (sobirát)}]{\ru{собирать} (sobirát) — to gather, to collect}

\label{lesson:RU-C209-sobirat}



\begin{quote}
Before the new one: say the Russian for to touch, then the Russian for to photograph.
\end{quote}

\subsection*{You'll want to know: \ru{собирать}}
\textbf{\ru{собирать}} (\emph{sobirát}) — \textquotedblleft{}to gather, to collect\textquotedblright{}.

\begin{grammarlens}[title={What you've built}]
One more word for this chapter.
\end{grammarlens}

\subsection*{Your turn}
\begin{itemize}
\raggedright
  \item \emph{Say it:} \emph{sobirát}
  \item \emph{Say it:} \emph{sobirát}, once more
  \item \emph{Say it:} say \emph{sobirát}
  \item \emph{Recall:} say the Russian for to touch, then the Russian for to photograph, then say \emph{sobirát} again
\end{itemize}

\begin{tcolorbox}[breakable,colback=teal!4,colframe=teal!35!black,title={Before you move on}]
What does \ru{собирать} mean? (\textquotedblleft{}To gather, to collect\textquotedblright{}.) Say it once more.
\end{tcolorbox}
\section[\texorpdfstring{\ru{обещать} (obeshchát)}{обещать (obeshchát)}]{\ru{обещать} (obeshchát) — to promise}

\label{lesson:RU-C209-obeshchat}



\begin{quote}
Before the new one: say the Russian for to photograph, then the Russian for to gather.
\end{quote}

\subsection*{You'll want to know: \ru{обещать}}
\textbf{\ru{обещать}} (\emph{obeshchát}) — \textquotedblleft{}to promise\textquotedblright{}.

\begin{grammarlens}[title={What you've built}]
One more word for this chapter.
\end{grammarlens}

\subsection*{Your turn}
\begin{itemize}
\raggedright
  \item \emph{Say it:} \emph{obeshchát}
  \item \emph{Say it:} \emph{obeshchát}, once more
  \item \emph{Say it:} say \emph{obeshchát}
  \item \emph{Recall:} say the Russian for to photograph, then the Russian for to gather, then say \emph{obeshchát} again
\end{itemize}

\begin{tcolorbox}[breakable,colback=teal!4,colframe=teal!35!black,title={Before you move on}]
What does \ru{обещать} mean? (\textquotedblleft{}To promise\textquotedblright{}.) Say it once more.
\end{tcolorbox}
\section[\texorpdfstring{\ru{шутить} (shutít)}{шутить (shutít)}]{\ru{шутить} (shutít) — to joke}

\label{lesson:RU-C209-shutit}



\begin{quote}
Before the new one: say the Russian for to gather, then the Russian for to promise.
\end{quote}

\subsection*{You'll want to know: \ru{шутить}}
\textbf{\ru{шутить}} (\emph{shutít}) — \textquotedblleft{}to joke\textquotedblright{}.

\begin{grammarlens}[title={What you've built}]
One more word for this chapter.
\end{grammarlens}

\subsection*{Your turn}
\begin{itemize}
\raggedright
  \item \emph{Say it:} \emph{shutít}
  \item \emph{Say it:} \emph{shutít}, once more
  \item \emph{Say it:} say \emph{shutít}
  \item \emph{Recall:} say the Russian for to gather, then the Russian for to promise, then say \emph{shutít} again
\end{itemize}

\begin{tcolorbox}[breakable,colback=teal!4,colframe=teal!35!black,title={Before you move on}]
What does \ru{шутить} mean? (\textquotedblleft{}To joke\textquotedblright{}.) Say it once more.
\end{tcolorbox}
\section[\texorpdfstring{\ru{жениться} (zhenítsya)}{жениться (zhenítsya)}]{\ru{жениться} (zhenítsya) — to get married (of a man)}

\label{lesson:RU-C209-zhenitsya}



\begin{quote}
Before the new one: say the Russian for to promise, then the Russian for to joke.
\end{quote}

\subsection*{You'll want to know: \ru{жениться}}
\textbf{\ru{жениться}} (\emph{zhenítsya}) — \textquotedblleft{}to get married (of a man)\textquotedblright{}.

\begin{grammarlens}[title={What you've built}]
One more word for this chapter.
\end{grammarlens}

\subsection*{Your turn}
\begin{itemize}
\raggedright
  \item \emph{Say it:} \emph{zhenítsya}
  \item \emph{Say it:} \emph{zhenítsya}, once more
  \item \emph{Say it:} say \emph{zhenítsya}
  \item \emph{Recall:} say the Russian for to promise, then the Russian for to joke, then say \emph{zhenítsya} again
\end{itemize}

\begin{tcolorbox}[breakable,colback=teal!4,colframe=teal!35!black,title={Before you move on}]
What does \ru{жениться} mean? (\textquotedblleft{}To get married (of a man)\textquotedblright{}.) Say it once more.
\end{tcolorbox}
